\section*{第 \ref{ch:g12:exp} 章 --- 指数与对数}

\begin{solution}{exo:g12:exp:1}
$\dfrac{\eu^{3x}\,\eu^{-x+1}}{\eu^{x}} = \eu^{3x - x + 1 - x} = \eu^{x+1}$.

$\ln\!\left(\dfrac{\eu^2\sqrt{\eu}}{\eu^{-3}}\right)
= 2 + \tfrac12 + 3 = \tfrac{11}{2}$.
\end{solution}

\begin{solution}{exo:g12:exp:2}
\emph{(a)} 设 $X = \eu^x > 0$，由 $X^2 - 3X + 2 = 0$ 得
$X = 1$ 或 $X = 2$，两者均为正，故 $x = 0$ 或 $x = \ln 2$。

\emph{(b)} 定义域为 $x - 1 > 0$ 且 $x + 2 > 0$，故 $x > 1$。方程
化为 $\ln\bigl((x-1)(x+2)\bigr) = \ln 4$，于是 $(x-1)(x+2) = 4$，
\emph{即} $x^2 + x - 6 = 0$，故 $x = 2$ 或 $x = -3$。仅 $x = 2$ 属于
定义域，因此 $S = \{2\}$。
\end{solution}

\begin{solution}{exo:g12:exp:3}
两边除以 $\eu^x$：
$\dfrac{1 - x^2\eu^{-x}}{1 + \eu^{-x}} \to \dfrac{1-0}{1+0} = 1$，利用
$x^2 \eu^{-x} \to 0$（见\cref{thm:g12:exp:growth}）。

$\dfrac{\ln x}{\sqrt x} \to 0$，由\cref{prop:g12:exp:lnanalytic}（
$n = 2$ 的情形）。

$x^2 \ln x = x \cdot (x \ln x) \to 0 \times 0 = 0$.

$x - \ln x = x\left(1 - \frac{\ln x}{x}\right) \to +\infty$，因为
$\frac{\ln x}{x} \to 0$。
\end{solution}

\begin{solution}{exo:g12:exp:4}
$f'(x) = \eu^{-x} - x\eu^{-x} = (1 - x)\eu^{-x}$，其符号与 $1 - x$ 一致：
$f$ 在 $\intoc{-\infty}{1}$ 上递增，在 $\intco{1}{+\infty}$ 上递减，
从而有全局最大值 $f(1) = \eu^{-1}$。

极限：当 $x \to +\infty$ 时，$f(x) = \frac{x}{\eu^x} \to 0$
（见\cref{thm:g12:exp:growth}）；当 $x \to -\infty$ 时，$x \to -\infty$ 且
$\eu^{-x} \to +\infty$，故 $f(x) \to -\infty$。

$f''(x) = -\eu^{-x} - (1-x)\eu^{-x} = (x - 2)\eu^{-x}$，它在
$x = 2$ 处变号，拐点为 $\bigl(2,\, 2\eu^{-2}\bigr)$。
\end{solution}

\begin{solution}{exo:g12:exp:5}
设 $g(x) = x - \ln(1+x)$，定义在 $\intoo{-1}{+\infty}$ 上。则
$g'(x) = 1 - \frac{1}{1+x} = \frac{x}{1+x}$，在 $\intoo{-1}{0}$ 上为负，且
在 $\intoo{0}{+\infty}$ 上为正：$g$ 取得最小值 $g(0) = 0$，
从而 $g \geq 0$，且等号仅在 $0$ 处成立。因此 $\ln(1+x) \leq x$。

取 $x = \frac1n$：
$n\ln\left(1 + \frac1n\right) \leq 1$，于是
$\left(1+\frac1n\right)^n = \eu^{\,n\ln(1+1/n)} \leq \eu$.
再取 $x = -\frac{1}{n+1} > -1$：
$\ln\left(1 - \frac{1}{n+1}\right) \leq -\frac{1}{n+1}$，于是
$-(n+1)\ln\left(1 - \frac{1}{n+1}\right) \geq 1$，从而
$\left(1 - \frac{1}{n+1}\right)^{-(n+1)} \geq \eu$.
\end{solution}

\begin{solution}{exo:g12:exp:6}
\emph{1.} 每年本金乘以 $1.03$：$C_n = C_0 (1.03)^n$。

\emph{2.} $C_n \geq 2C_0 \iff (1.03)^n \geq 2 \iff n \ln 1.03 \geq \ln 2
\iff n \geq \frac{\ln 2}{\ln 1.03} \approx 23.45$：本金在
$24$ 年后翻倍。

\emph{3.} $\frac{70}{3} \approx 23.3$，与精确值
$\frac{\ln 2}{\ln 1.03}$。该法则成立的原因是
$\ln(1 + r) \approx r$（$r$ 较小时），故
$\frac{\ln 2}{\ln(1+r)} \approx \frac{0.693}{r} \approx \frac{70}{100r}$.
\end{solution}

\begin{solution}{exo:g12:exp:7}
$\eu^{2x} > \eu^{x+1} \iff 2x > x + 1 \iff x > 1$（指数函数
严格递增）：$S = \intoo{1}{+\infty}$。

第二个不等式的定义域为 $x^2 - 1 > 0$ 且 $x + 5 > 0$，\emph{即}
$x \in \intoo{-5}{-1} \cup \intoo{1}{+\infty}$。在此定义域上，$\ln$
严格递增，于是
\[
\ln(x^2-1) \leq \ln(x+5) \iff x^2 - 1 \leq x + 5 \iff x^2 - x - 6 \leq 0
\iff x \in \intcc{-2}{3}.
\]
与定义域取交集，得 $S = \intco{-2}{-1} \cup \intoc{1}{3}$。
\end{solution}

\begin{solution}{exo:g12:exp:8}
\emph{1.} $f'(x) = \dfrac{\eu^x x - \eu^x}{x^2} = \dfrac{(x-1)\eu^x}{x^2}$，
在 $\intoo{0}{1}$ 上为负，在 $\intoo{1}{+\infty}$ 上为正，故有最小值
$f(1) = \eu$。极限：$f \to +\infty$ 于 $0^+$ 处（分子 $\to 1$，
分母 $\to 0^+$），在 $+\infty$ 处亦然（见\cref{thm:g12:exp:growth}）。

\emph{2.} 当 $x > 0$ 时，$\eu^x = kx \iff f(x) = k$。由变化表
（$+\infty \searrow \eu \nearrow +\infty$，两侧各自连续且严格单调）可知：
$k < \eu$ 时无解；恰有一解 $x = 1$，此时 $k = \eu$；
$k > \eu$ 时恰有两解（一解在 $\intoo{0}{1}$ 中，
另一解在 $\intoo{1}{+\infty}$ 中，由每个区间上的双射定理得）。
\end{solution}

\begin{solution}{exo:g12:exp:9}
\emph{1.} 由\cref{exo:g12:exp:5}中的第一个不等式直接可得：
$u_n = \eu^{\,n\ln(1+1/n)} \leq \eu^1 = \eu$.

\emph{2.} 写作
\[
\ln u_n = n \ln\!\left(1 + \frac1n\right)
= \frac{\ln\!\left(1 + \frac1n\right) - \ln 1}{\frac1n},
\]
这是 $\ln$ 在点 $1$ 处的差商，其增量为
$h = \frac1n \to 0$。因 $\ln$ 在 $1$ 处可导，且导数为
$1$，$\ln u_n \to 1$。由 $\exp$ 的连续性
（见\cref{prop:g12:limcont:seqcont}），$u_n = \eu^{\ln u_n} \to \eu^1 = \eu$。
\end{solution}

\begin{solution}{pb:g12:exp:1}
\textbf{1.} $x = \frac{\ln 5}{2}$。\quad
$3x - 1 = \eu^2$：$x = \frac{\eu^2 + 1}{3}$。\quad 令
$u = \eu^x$：$u^2 - 3u + 2 = (u - 1)(u - 2) = 0$，故 $x = 0$ 或
$x = \ln 2$。

\textbf{2.} $\frac{3\ln 2}{\ln 2} = 3$；\quad
$3 + \frac12 = \frac72$；\quad $\frac52$。

\textbf{3.} $f'(x) = \eu^x - 1$：在 $0$ 之前为负，
之后为正，极小值为 $f(0) = 0$，故 $f \geq 0$ 处处成立，即
$\eu^x \geq 1 + x$。代入 $x = \ln(1 + u)$：
$1 + u \geq 1 + \ln(1 + u)$，即 $\ln(1 + u) \leq u$。

\textbf{4.} $x^2\eu^{-x} = \frac{x^2}{\eu^x} \to 0$，且
$\frac{\ln x}{x} \to 0$：指数压倒幂，幂
压倒对数。又 $\frac{\eu^x - 1}{x} =
\frac{\eu^x - \eu^0}{x - 0} \to \exp'(0) = 1$。

\textbf{5.} $f'(x) = \ln x + 1$：零点为 $\eu^{-1}$，之前为负，
之后为正，极小值为
$f\!\left(\frac1\eu\right) = -\frac1\eu$；且
$x \ln x \to 0$（在 $0^+$ 处）：曲线从原点出发，下降至
$-\frac1\eu$，随后上升远去。

\textbf{6.} $n = \frac{\ln 2}{\ln 1.03} \approx 23.4$ 年。

\textbf{7.} 精确值：$23.4$、$11.9$、$8.0$ 年；72 法则：
$24$、$12$、$8$。由于 $\ln(1+r) \approx r$，
$\frac{\ln 2}{\ln(1+r)} \approx \frac{0.693}{r}$，即\
$\frac{69.3}{\text{利率 }\%}$；银行家则向上取整为 $72$，
因为它能很好地被 $2, 3, 4, 6, 8, 9, 12$ 整除 ---
心算胜过小数的精确度。

\textbf{8.} $2^{t/20} = \eu^{t \ln 2 / 20}$：
每分钟 $\lambda = \frac{\ln 2}{20}$。经过
$24 \times 60 = 1440$ 分钟后：$2^{72} \approx 4.7 \times
10^{21}$ 个细菌 --- 比地球上的沙粒还多，
由一个细胞在一天之内繁衍而来。该模型只有在食物
和空间充足时才可靠：真实增长会弯曲为逻辑斯谛 S 形曲线
（即\cref{pb:g12:deriv:1}中流行病学家的曲线）。

\textbf{9.} 在 23:00（经过 $8$ 小时）：
$100 \times 2^{-8/5} \approx 33$ mg。低于 $10$ mg 时：
$2^{-t/5} < 0.1$ 给出 $t > 5\,\frac{\ln 10}{\ln 2} \approx
16.6$ 小时，即次日早上 7:40 左右 --- 下午
三点的那杯浓缩咖啡，余韵悠长。

\textbf{10.} 按年计息：$2$。按月计息：
$\left(1 + \frac{1}{12}\right)^{12} \approx 2.613$。按日计息：
$\approx 2.715$。复利的上限是 $\eu = 2.71828\ldots$
（见\cref{exo:g12:exp:9}） --- 连续复利，贪婪
也就收敛了。

\textbf{11.} 若病例数按指数增长，$c(t) = c_0
\eu^{\lambda t}$，则 $\ln c(t) = \ln c_0 + \lambda t$：
一条斜率为 $\lambda$ 的直线 --- 即增长率。
对数图上的一段直线\emph{正是}指数增长
阶段，其陡峭程度就是疫情的节奏。

\textbf{12.} $120 - 60 = 60$ dB 意味着
$10\log_{10}(I_2/I_1) = 60$：音乐会的声强
是交谈的 $10^6$ 倍 --- 即一百万倍 ---
声强远胜交谈。

\textbf{13.} 振幅之比：$10^{7-5} = 100$。能量之比：
$100^{3/2} = 1000$：两个震级点
在能量上隐藏着三个数量级。

\textbf{14.} $10^{6.5 - 2} = 10^{4.5} \approx 32\,000$：柠檬
汁的酸性是牛奶的三万倍 --- pH
把化学中的巨大鸿沟压缩进一把袖珍标尺。

\textbf{15.} $\log_2\!\frac{4186}{27.5} =
\frac{\ln(4186/27.5)}{\ln 2} \approx 7.25$：一架钢琴的音域刚刚
超过七个八度。感官对刺激的\emph{比值}作出反应
--- 强度加倍，感觉上只是一步，无论
起点多高 --- 因此感知建立在对数
尺度之上，我们为此发明的单位亦然。

\textbf{16.} 将代表 $a$ 的刻度尺段与
代表 $b$ 的刻度尺段首尾相接，长度相加为 $\ln a + \ln b$，
于是该总长度处的刻度读数为
$\eu^{\ln a + \ln b} = ab$：对数尺
通过把对数相加来计算乘积 --- $\ln(ab) = \ln a + \ln b$
（见\cref{prop:g12:exp:lnalg}）刻在黄杨木上。

\textbf{17.} $\frac{\eu^x}{x}$ 在
$\intoo{0}{+\infty}$ 上的最小值为 $\eu$，在 $x = 1$ 处取得：
当 $k < \eu$ 时无解，$k = \eu$ 时恰有一解，$k > \eu$ 时有两解。
切线检验：在 $a = 1$ 处，$\eu^x$ 的切线为
$y = \eu + \eu(x - 1) = \eu x$：它经过原点
--- 这条临界直线恰好擦过曲线上的点 $(1, \eu)$。

\textbf{18.} $n \ln\!\left(1 - \frac1n\right) =
\frac{\ln(1 - 1/n)}{-1/n} \times (-1) \to -1$，从而
$\left(1 - \frac1n\right)^n \to \eu^{-1} \approx 0.368$：这正是
\cref{pb:g11:binom:1} 中的 $0.999^{1000} \approx 0.368$，
得到了解释 --- 彩票“差一点中奖”的常数为
$\frac1\eu$。

\textbf{19.} 二者都是“许多独立
且各自不太可能发生的事件”的极限形态：
\emph{没有一次}触发的概率，在 $n$ 个大小为 $\frac1n$ 的机会中，趋于
$\eu^{-1}$，而秘书问题的拒绝窗口的设定
使成功集中于同一常数。
$\frac1\eu \approx 0.368$。

\textbf{20.} 复利的上限；唯一的底数，
其在 $0$ 处的切线斜率为 $1$（这正是微积分偏爱它的原因）；
超越一切幂的增长；“差一点中奖”与最优停止
问题共同归宿的常数；以及对数 ---
乘法化为加法，世界上最悬殊的量程
被折叠到一把尺子上。
\end{solution}
